\selectlanguage{ngerman}%
\frenchspacing%
\begin{ledgroupsized}[r]{120mm}%
\footnotesize%
\pstart%
\noindent\textbf{Überlieferung:}
\pend%
\end{ledgroupsized}%
\begin{ledgroupsized}[r]{114mm}%
\footnotesize%
\pstart%
\parindent -6mm%
\makebox[6mm][l]{\textit{L}}%
Aufzeichnung:
LH XXXV 10,\;14 Bl.\;11.
Ein Blatt 8\textsuperscript{o},
Ränder z.\,T. beschnitten und abgerissen.
Zwei voll beschriebene Seiten.
Bl.\;\href{\Scan{LH_35_10_14_011v}}{11\:v\textsuperscript{o}} überliefert ferner, z.\,T. überschrieben, die Anrede eines Briefes an D.\;Nessel:\indexN{\textls{Nessel} (Nesselius), Daniel 1644\textendash1700}
\textit{Amplissimo Dno Nesselio Consiliario et Bibliothecario Caesareo}. %
Im Kopf von Bl.\;\href{\Scan{LH_35_10_14_011r}}{11\,r\textsuperscript{o}} ist von unklarer Hand die römische Zahl \textit{X} in roter Tinte vermerkt
(vgl. die Träger der benachbarten Signaturen LH~XXXV~10,\;14 Bl.\;2 bis~13,
u.\,a. 
N.\;\ref{60118},
\ref{60115},
\ref{60112},
\ref{60116},
\ref{60113},
\ref{60114},
\ref{60110},
\ref{41824}
u.\;\ref{60111}
im vorliegenden Band).
\pend%
\end{ledgroupsized}%
\vspace{5mm}%
\begin{ledgroup}%
\footnotesize%
\pstart%
\noindent%
\textbf{Datierungsgründe:}
Die Datierung der Aufzeichnung N.\;\ref{60117} beruht auf Eigenschaften des Textträgers.
Siehe die ausführliche Begründung auf S.\;\pageref{60117_dat}.%
\pend%
\end{ledgroup}%
\selectlanguage{latin}%
\frenchspacing%
\vspace{5mm}%
\normalsize%
\pstart%
\noindent%
\lbrack%
\href{\Scan{LH_35_10_14_011r}}%
{11\:r\textsuperscript{o}}%
\rbrack\
Sint duo Globi
\textit{A} et \textit{B}
inaequales connexi%
\indexS{globi connexi filo}
\edtext{%
filo
ipsis diductis
tenso,%
\indexS{filum tensum}
tum subito dimittantur,%
}{%
\lemma{filo}%
\Bfootnote{\hspace{-0,5mm}%
\bars{ipsis diductis \textit{erg.}} tenso, %
\textit{(1)}~quod %
\textit{(2)}~tum subito %
\textit{(a)}~demittatur, %
\textit{(b)}~dimittantur,%
~\textit{L}%
}}
quaeritur
qua proportione celeritatum%
\indexS{proportio celeritatum}
ad se invicem accedant.
Suppono autem
tensiones esse%
\indexS{tensiones ut longitudines}
ut longitudines.
Ajo
\edtext{%
cum in conatibus primis%
\indexS{conatus primus}
vires sint reciproce
ut corpora%
\indexS{vires reciproce ut corpora}%
}{%
\lemma{cum}%
\Bfootnote{%
\textit{(1)}~conatus primi sint %
\textit{(a)}~ut %
\textit{(b)}~reciproce ut corpora %
\textit{(2)}~in conatibus %
\lbrack...\rbrack\ ut corpora%
~\textit{L}%
}}
et in quovis
\edtext{%
addititio conatu%
\indexS{conatus addititius}
ex nova contractione%
\indexS{contractio fili tensi}
similiter,%
}{%
\lemma{addititio}%
\Bfootnote{%
\textit{(1)}~similiter %
\textit{(a)}~et nova contra %
\textit{(b)}~con %
\textit{(2)}~conatu ex nova %
\textit{(a)}~attractione\indexS{attractio} %
\textit{(b)}~contractione similiter,%
~\textit{L}%
}}
hinc velocitates%
\indexS{velocitas acquisita}
\edtext{%
acquisitas%
}{%
\lemma{acquisitas}%
\Bfootnote{%
\textit{erg.\;L}%
}}
semper esse corporibus reciproce proportionales,
adeoque concursum fore in centro gravitatis.%
\indexS{concursus in centro gravitatis}%
\indexS{centrum gravitatis concursus}
Et idem
\edtext{%
fi\lbrack e\rbrack t,%
}{%
\lemma{fiat}%
\Bfootnote{%
\textit{L\;ändert Hrsg.}%
}}
si reflexione facta%
\indexS{reflexio post concursum}
iterum a se digrediantur
et filum denuo%
\indexS{filum tensum}
\edlabel{LH_35_10_14_011_011r_ujf-1}%
tendant.%
\edtext{}{%
\xxref{LH_35_10_14_011_011r_ujf-1}{LH_35_10_14_011_011r_ujf-2}%
\lemma{tendant.}%
\Bfootnote{\hspace{-0,5mm}%
\lbrack/\rbrack\ %
\textit{(1)}~Si filum in %
\textit{(2)}~Si corpus \textit{A} %
\textit{(a)}~certa %
\textit{(b)}~velocitate aliqua impellatur%
~\textit{L}%
}}
\pend%
\begin{edfigure}[0.35em][0.90em]
\includegraphics%
{gesamttex/edit_VIII-4/images/LH_35_10_14_011_d1_011r.pdf}
\edcaption[0.00em]{%
\hspace{2,75mm}%
\lbrack\textit{Fig.\:1}\rbrack%
}%
\end{edfigure}%
\pstart%
\noindent%
Si corpus \textit{A}
velocitate aliqua%
\indexS{velocitas impulsus}
impellatur%
\edlabel{LH_35_10_14_011_011r_ujf-2}
et ope
\edtext{%
fili \textit{AB}%
\indexS{filum tensum}
corpus \textit{B}%
}{%
\lemma{fili}%
\Bfootnote{%
\textit{(1)}~tensi corpus %
\textit{(2)}~tendibilis %
\textit{(3)}~\textit{AB} corpus \textit{B}%
~\textit{L}%
}}
secum trahat,%
\indexS{tractio corporis ab altero}%
\indexS{corpus tractum ab altero}
patet,
si filum rigidum esset,%
\indexS{filum rigidum}
motum eorum fore talem,%
\indexS{motus corporum connexorum}
ut vires sint%
\indexS{vires aequales}
\edtext{%
aequales\lbrack;\rbrack\ %
adeoque potestas%
\indexS{potestas movendi}
qua
move\lbrack re\rbrack tur
$A + B$
\lbrack foret\rbrack\
ad eam%
}{%
\lemma{aequales,}%
\Bfootnote{%
\textit{(1)}~ut aute %
\textit{(2)}~adeoque %
\textit{(a)}~potestatem %
\textit{(b)}~potestates reciproce %
\textit{(c)}~potestas qua %
\textit{(aa)}~moveretur %
\textit{(bb)}~%
\bars{movetur \textit{ändert Hrsg.}} $A + B$ %
\bars{foret \textit{erg. Hrsg.}} ad eam%
~\textit{L}%
}}
qua \textit{A} moveretur,
ut \textit{A} ad $A + B$ %
seu reciproce ut corpora.
Nunc vero
cum filum sit tendibile%
\indexS{filum tendibile}
et partem virium%
\indexS{vis tendens filum}%
\indexS{vis recepta in filum tendibile}
in se
\edtext{%
recipiat%
\lbrack,\rbrack\
magis%
}{%
\lemma{recipiat}%
\Bfootnote{%
\textit{(1)}~minus %
\textit{(2)}~magis%
~\textit{L}%
}}
movebitur \textit{A}
quam in casu rigoris,%
\indexS{rigor}
et \textit{B} minus.
\ledpb%
Quaeritur
quomodo distribuendae vires.%
\indexS{vis distribuenda}%
\indexS{distributio virium}
Ponamus aliquousque extendi filum%
\indexS{filum extensum}
motu ipsius \textit{A},
nec tamen trahi \textit{B},
quippe retentum,
ita ut sola vis
\edtext{%
\lbrack ipsius \textit{A}\rbrack%
}{%
\lemma{ipsius \textit{A}}%
\Bfootnote{%
\textit{erg.\;Hrsg.}%
}}
in filum perveniat,%
\indexS{filum extensum}
\edtext{%
quae vis%
\indexS{vis tendens filum}%
}{%
\lemma{quae}%
\Bfootnote{%
\textit{(1)}~si %
\textit{(2)}~vis%
~\textit{L}%
}}
tantum valeat %
\edtext{%
quantum $at^2$%
}{%
\lemma{quantum}%
\Bfootnote{%
\textit{(1)}~$av^2$ %
\textit{(2)}~$at^2$%
~\textit{L}%
}}
seu quantum celeritas \textit{t}
\edtext{%
ipsius
\lbrack\textit{A};\rbrack\
hoc jam posito%
}{%
\lemma{ipsius}%
\Bfootnote{%
\bars{\textit{t}, \textit{ändert Hrsg.}} hoc
\textit{(1)}~pacto %
\textit{(2)}~jam posito%
~\textit{L}%
}}
dimittatur,
patet
filum vi %
$t^2a$%
\indexS{vis tendens filum}%
\indexS{filum contrahens corpora}
\edtext{%
contrahere
corpora~\textit{A}
et~\textit{B},%
}{%
\lemma{contrahere}%
\Bfootnote{%
\bars{reciproce \textit{gestr.}} corpora~\textit{A} et~\textit{B},%
~\textit{L}%
}}
et corpori quidem \textit{A}
imprimere velocitatem%
\indexS{velocitas impressa}
\edtext{%
versus \textit{B} %
quae sit ut \textit{b},%
}{%
\lemma{versus \textit{B},}%
\Bfootnote{%
\textit{(1)}~ut \textit{a} %
\textit{(2)}~quae sit ut %
\textit{(a)}~\textit{a} %
\textit{(b)}~\textit{b},%
~\textit{L}%
}}
et corpori \textit{B}
velocitatem%
\indexS{velocitas impressa}
versus
\edtext{%
\lbrack\textit{A}\rbrack%
}{%
\lemma{\textit{B}}%
\Bfootnote{%
\textit{L\;ändert Hrsg.}%
}}
quae sit ut \textit{a}.
Verum
cum corpus \textit{A} interim pergat%
\indexS{corpus pergens}%
\lbrack,\rbrack\
secum abripiet corpus \textit{B},%
\indexS{corpus abripiens alterum}%
\indexS{corpus abreptum ab altero}
ita tamen
ut minor sit velocitas ipsius \textit{A} %
quam ante%
\lbrack,\rbrack\
et tamen minor ipsius \textit{B} %
quam ipsius \textit{C}.
Et quidem missis
\edlabel{LH_35_10_14_011_011r_nmf-1}%
praedictis,
rem ita concipiamus.%
\edtext{}{%
\xxref{LH_35_10_14_011_011r_nmf-1}{LH_35_10_14_011_011r_nmf-2}%
\lemma{praedictis,}%
\Bfootnote{%
\textit{(1)}~cum aeque %
\textit{(2)}~rem ita concipiamus. %
\lbrack/\rbrack\ %
\textit{(a)}~Ut \textlangle vi\textrangle\ corpus %
\textit{(b)}~Ponamus %
\textit{(aa)}~\textit{AB} esse rigidum sed fragile, %
\textit{(bb)}~rectam \textit{AB} %
\lbrack...\rbrack\ sed fragilem, %
\textit{(aaa)}~ut %
\textit{(bbb)}~ut %
\textit{(ccc)}~ita ut%
~\textit{L}%
}}%
\pend%
\begin{edfigure}[1.5em][1.00em]
\includegraphics%
{gesamttex/edit_VIII-4/images/LH_35_10_14_011_d2_011r.pdf}
\edcaption[0.25em]{%
\lbrack\textit{Fig.\:2}\rbrack%
}%
\end{edfigure}%
\pstart%
\noindent%
Ponamus rectam \textit{AB} esse rigidam%
\indexS{recta rigida}
sed fragilem,%
\indexS{recta fragilis}
ita ut%
\edlabel{LH_35_10_14_011_011r_nmf-2}
percusso fortiter \textit{A}
celeritate quadam certa,%
\indexS{celeritas percussionis}
rumpi possit %
ob resistentiam ad motum%
\indexS{resistentia ad motum}
seu inertiam naturalem%
\indexS{inertia naturalis}
ipsius~\textit{B},
\edtext{%
ajo
si percutiatur \textit{A}
directione%
\indexS{directio percussionis}
\lbrack\textit{BA}\rbrack\ %
vi tali%
\indexS{vis percussionis}
qua frangatur \textit{AB},
impressum iri toti \textit{AB}
vim%
}{%
\lemma{ajo}%
\Bfootnote{%
\textit{(1)}~tantum %
\textit{(2)}~tot %
\textit{(3)}~si %
\textit{(a)}~ictu sufficiente %
\textit{(b)}~\bars{si \textit{streicht Hrsg.}} percutiatur \textit{A} directione %
\bars{\textit{BA}, \textit{gestr., wieder gültig macht Hrsg.}}
\textit{(aa)}~im %
\textit{(bb)}~ea %
\textit{(cc)}~imp %
\textit{(dd)}~vi tali qua frangatur \textit{AB}, %
\textit{(aaa)}~impressam %
\textit{(bbb)}~impressum iri 
\textit{(aaaa)}~ipsi \textit{B} %
\textit{(bbbb)}~toti \textit{AB} vim%
~\textit{L}%
}}
quae impressa fuisset%
\indexS{vis impressa}
si \textit{AB} non fuisset fragilis,%
\indexS{recta fragilis}
demta vi%
\indexS{vis demta}
impensa in fracturam%
\indexS{fractura}%
\indexS{vis impensa in fracturam}%
\lbrack;\rbrack\
vis autem impensa in fracturam%
\indexS{fractura}%
\indexS{vis impensa in fracturam}
soli \textit{B} detrahetur,%
\indexS{vis detracta}
et
\edtext{%
quidem differentia celeritatis%
\indexS{differentia celeritatis}%
\indexS{celeritas impressa}
inter~\textit{A} et~\textit{B}
erit celeritas ipsius~\textit{B}%
}{%
\lemma{quidem}%
\Bfootnote{%
\textit{(1)}~tanta celerit %
\textit{(2)}~eri %
\textit{(3)}~differentia celeritatis %
\lbrack...\rbrack\ ipsius \textit{B}%
~\textit{L}%
}}
quae sufficeret%
\indexS{celeritas sufficiens ad frangendum}
\edtext{%
ad fra\lbrack n\rbrack gendam
lineam \textit{AB}.%
\indexS{linea frangenda}
Si jam
continuato vinculo rigido%
\indexS{vinculum rigidum}
\textit{AB} in \textit{BC},
etiam corpus \textit{C} eodem ictu moveri%
\indexS{corpora mota eodem ictu}%
\indexS{ictus}
et linea tam \textit{AB} %
quam \textit{BC} frangi%
\indexS{linea fracta}
intelligatur,
patet
ipsi \textit{A} nihil detrahi de celeritate,%
\indexS{celeritas detracta}
ipsi \textit{B} aeque ac ipsi \textit{C} detrahi vim%
\indexS{vis detracta}%
\indexS{vis frangendi}
qua frangitur \textit{AB},
at soli \textit{C} detrahi vim%
\indexS{vis detracta}%
\indexS{vis frangendi}
qua frangitur \textit{BC};%
}{%
\lemma{ad}%
\Bfootnote{%
\bars{fragendam \textit{ändert Hrsg.}} lineam \textit{AB}. %
\textit{(1)}~Assumto jam tertio corpore \textit{B}, ponamus %
\textit{(2)}~Si jam continuato vinculo rigido %
\textit{(a)}~etiam %
\textit{(b)}~\textit{AB} in %
\lbrack...\rbrack\ moveri et %
\textit{(aa)}~frangi %
\textit{(bb)}~linea tam %
\lbrack...\rbrack\ intelligatur, patet %
\textit{(aaa)}~corpori %
\textit{(aaaa)}~vi %
\textit{(bbbb)}~motum %
\textit{(cccc)}~vim qua frangitur \textit{B} %
\textit{(bbb)}~ipsi \textit{A} %
\lbrack...\rbrack\ de celeritate, %
\textit{(aaaa)}~at %
\textit{(bbbb)}~ult %
\textit{(cccc)}~ipsi \textit{B} %
\lbrack...\rbrack\ at soli %
\textit{(aaaaa)}~\textit{B} %
\textit{(bbbbb)}~\textit{C} detrahi vim qua frangitur \textit{BC};%
~\textit{L}%
}}
et eadem foret ratiocinatio,%
\indexS{ratiocinatio}
si plura corpora essent juncta.%
\indexS{corpora juncta}
Porro vis%
\indexS{vis frangendi}
qua frangitur \textit{BC}
sit eadem cum ea
qua frangitur%
\indexS{vis frangendi}
\edtext{%
\textit{AB},
itaque tota vis frangendo perdita%
\indexS{vis perdita frangendo}
sit \textit{v},
\ledpb%
ejus dimidium%
\indexS{dimidium virium}
reservetur ipsi \textit{C},
de reliquo dimidio
\lbrack%
\href{\Scan{LH_35_10_14_011v}}%
{11\:v\textsuperscript{o}}%
\rbrack\
etiam dimidium%
\indexS{dimidium virium}
ipsi \textit{C} assignetur,
\rule[-0mm]{0pt}{5mm}%
et unum tantum ipsi \textit{B},%
}{%
\lemma{\textit{AB},}%
\Bfootnote{%
\textit{(1)}~patet illam qua frangitur \textit{AB} %
\textit{(a)}~dividi in %
\textit{(b)}~diminu %
\textit{(2)}~itaque tota %
\lbrack...\rbrack\ reservetur ipsi %
\textit{(a)}~\textit{BC} %
\textit{(b)}~\textit{C}, de %
\lbrack...\rbrack\ ipsi \textit{B},%
~\textit{L}%
}}
et
\edtext{%
proinde vis detrahenda%
\indexS{vis detrahenda}%
}{%
\lemma{proinde}%
\Bfootnote{%
\textit{(1)}~celeritas d %
\textit{(2)}~vis detrahenda%
~\textit{L}%
}}
ipsi \textit{B} erit
$\displaystyle\frac{1}{4}$,
detrahenda ipsi \textit{C} erit
\rule[-0mm]{0pt}{5mm}%
$\displaystyle \frac{1}{4} + \displaystyle\frac{2}{4}$.
Si
\edtext{%
vero sint quatuor%
}{%
\lemma{vero}%
\Bfootnote{%
\textit{(1)}~tria %
\textit{(2)}~sint quatuor%
~\textit{L}%
}}
corpora
\edtext{%
[\textit{A},]%
}{%
\lemma{\textit{A},}%
\Bfootnote{\textit{erg.\;Hrsg.}%
}}
\textit{B}, \textit{C}, \textit{D},
tunc soli \textit{D}
servetur
\edtext{%
$\displaystyle\frac{1}{2}$,
ipsis \textit{D} et \textit{C} singulis%
}{%
\lemma{$\displaystyle\frac{1}{2}$,}%
\Bfootnote{%
\textit{(1)}~vel %
\textit{(2)}~soli \textit{C} %
\textit{(3)}~ipsis \textit{D} et \textit{C} singulis%
~\textit{L}%
}}
$\displaystyle\frac{1}{4}$,
ipsis \textit{D}, \textit{C}, \textit{B}
singulis $\displaystyle\frac{1}{8}$.
Ergo ipsi \textit{D} detrahetur vis
$\displaystyle \frac{1}{2} + \displaystyle\frac{1}{4} + \displaystyle\frac{1}{8}$,
ipsi \textit{C}
$\displaystyle \frac{1}{4} + \displaystyle\frac{1}{8}$,
et denique ipsi \textit{B} %
$\displaystyle\frac{1}{8}$ de tota vi %
quae erat et ad priorem%
\lbrack,\rbrack\
cum duo solum essent corpora,
\rule[-0mm]{0pt}{4.5mm}%
ut 3 ad 2.
\newline%
\indent%
\rule[-0mm]{0pt}{9.0mm}%
Sint corpora infinita%
\indexS{corpora infinita} %
\begin{tabular}{ccccccccccc}
&$\displaystyle\frac{1}{2}$&&$\displaystyle\frac{1}{4}$&&$\displaystyle\frac{1}{8}$&&$\displaystyle\frac{1}{16}$&&&\\
$\displaystyle\frac{1}{2}$&&$\displaystyle\frac{1}{4}$&&$\displaystyle\frac{1}{8}$&&$\displaystyle\frac{1}{16}$&&$\displaystyle\frac{1}{32}$&&$\displaystyle\frac{1}{64}$
\end{tabular}
etc.%
\\
\rule[-0mm]{0pt}{5mm}%
$1=$
\rule[0.5mm]{68.5mm}{0.6pt}\
etc. vis detracta ipsi ultimo ex infinitis%
\\
\rule[-0mm]{0pt}{8mm}\rule[0.5mm]{9mm}{0pt}%
$\displaystyle\frac{1}{2}=$
\rule[0.5mm]{58.5mm}{0.6pt}\
etc. penultimo%
\\
\rule[-0mm]{0pt}{6mm}\rule[0.5mm]{18mm}{0pt}%
$\displaystyle\frac{1}{4}=$
\rule[0.5mm]{49.5mm}{0.6pt}\
etc. antepenultimo etc.%
\newline%
Et ut verbo dicam,%
\indexS{verbum}
\rule[-0mm]{0pt}{5mm}%
\edtext{%
si numeri inde a primo sumantur
in progressione arithmetica,%
\indexS{progressio arithmetica}%
}{%
\lemma{si}%
\Bfootnote{%
\textit{(1)}~corpora sint u %
\textit{(2)}~numeri %
\textit{(a)}~a %
\textit{(b)}~inde a %
\lbrack...\rbrack\ progressione arithmetica,%
~\textit{L}%
}}
erunt vires detractae
in progressione%
\indexS{progressio geometrica}
\edtext{%
Geometrica,
et proinde si corpora sint aequalia,
erunt celeritatum detractiones%
\indexS{detractio celeritatis}%
\indexS{celeritas detracta}
ut numeri,%
}{%
\lemma{Geometrica}%
\Bfootnote{%
\textit{(1)}~. Sed quid si ponamus %
\textit{(2)}~, et proinde %
\lbrack...\rbrack\ aequalia, erunt %
\textit{(a)}~differentiae %
\textit{(b)}~celeritatum %
\textit{(aa)}~a celeri %
\textit{(bb)}~detractiones %
\textit{(aaa)}~in progressi %
\textit{(bbb)}~ut %
\textit{(ccc)}~ut numeri,%
~\textit{L}%
}}
si distantiae a primo corpore
sint ut
\edtext{%
logarithmi.%
\indexS{logarithmus}
\newline%
\indent%
\textls{Imo erratum videtur.}%
\indexS{erratum}%
}{%
\lemma{logarithmi.}%
\Bfootnote{%
\textit{(1)}~Si filum %
\textit{(2)}~Si fila ponantur esse multiplicata, in ratione distantiarum, %
\textit{(a)}~ita ut corpus %
\textit{(b)}~tunc toti corpori ultimo %
\textit{(c)}~tunc numerus corporum sit \textit{z}, et fila erunt \textit{z}, $z-1$, $z-2$, $z-3$, etc. tota igitur vis erit %
\lbrack/\rbrack\
\textit{(3)}~\textls{Imo erratum videtur.}%
~\textit{L}%
}}
Si sint duo corpora%
\indexS{corpus trahendum}%
\indexS{tractio corporum}
\edtext{%
trahenda a primo \textit{A},
nempe \textit{B} et \textit{C},
vis%
}{%
\lemma{trahenda}%
\Bfootnote{%
\textit{(1)}~, \textit{B} et \textit{C}, vis %
\textit{(2)}~a primo %
\lbrack...\rbrack\ \textit{C}, vis%
~\textit{L}%
}}
utique dividenda%
\indexS{vis dividenda}%
\indexS{vis tractionis}
in duas partes aequales
et ipsi \textit{C} detrahetur%
\indexS{vis detracta}
\rule[-0mm]{0pt}{5.5mm}%
$\displaystyle \frac{1}{2} + \displaystyle\frac{1}{4}$
et ipsi \textit{B}
$\displaystyle \frac{1}{4}$.
Sed si sint tria trahenda%
\indexS{corpus trahendum}%
\indexS{tractio corporum}
\textit{B}, \textit{C}, \textit{D},
vis dividenda in partes tres,%
\indexS{vis dividenda}%
\indexS{vis tractionis}
\edtext{%
et soli ultimo \textit{D} servabitur%
}{%
\lemma{et}%
\Bfootnote{%
\textit{(1)}~uno %
\textit{(2)}~soli ultimo %
\bars{\textit{D} \textit{erg.}} servabitur%
~\textit{L}%
}}
\rule[-0mm]{0pt}{5mm}%
$\displaystyle\frac{1}{3}$,
ipsis \textit{D} et \textit{C} servabitur
\rule[-0mm]{0pt}{5mm}%
$\displaystyle\frac{1}{3}$,
ipsis \textit{D}, \textit{C}, \textit{B},
$\displaystyle\frac{1}{3}$,
\ledpb%
ergo \textit{D} accipiet
$\displaystyle \frac{1}{3} + \frac{1}{6} + \frac{1}{9}$,
et \textit{C}
$\displaystyle \frac{1}{6} + \displaystyle\frac{1}{9}$,
et \textit{B}~%
$\displaystyle\frac{1}{9}$.
Ergo
si numerus sit \textit{z},
erit vis detrahenda%
\indexS{vis detrahenda}
ipsi \textit{z} %
ut
\rule[-0mm]{0pt}{5.5mm}%
$\displaystyle \frac{1}{z} + \frac{1}{z-1} + \frac{1}{z-2}$
etc.
Erit ergo area%
\indexS{area hyperbolae}
\edtext{%
Hyperbolae%
\indexS{hyperbola}
repr\lbrack aesent\rbrack ans%
}{%
\lemma{Hyperbolae}%
\Bfootnote{%
\textit{(1)}~distribuenda %
\textit{(2)}~\bars{reprehensans \textit{ändert Hrsg.}}%
~\textit{L}%
}}
totam vim%
\indexS{vis tota}%
\indexS{area repraesentans vim}%
\lbrack;\rbrack\
\rule[-0mm]{0pt}{3.5mm}%
et areae %
inde ab ultima ordinata
\edtext{%
sumtae,
ita ut distantia extremarum ordinatarum
\rule[-0mm]{0pt}{3.5mm}%
sit distantia corporis a primo,
dab\lbrack un\rbrack t vim%
}{%
\lemma{sumtae}%
\Bfootnote{%
\textit{(1)}~cresc %
\textit{(2)}~abscissa %
\textit{(3)}~\textlangle axi\textrangle\ %
\textit{(4)}~q %
\textit{(5)}~ita ut %
\lbrack...\rbrack\ ordinatarum sit %
\textit{(a)}~numerus cor %
\textit{(b)}~distantia corporis a primo, %
\bars{dabit \textit{ändert Hrsg.}} vim%
~\textit{L}%
}}
ei detrahendam corpori.%
\indexS{vis detrahenda}%
\rule[-0mm]{0pt}{3.5mm}%
\pend%
\begin{edfigure}[1.25em][1.00em]
\includegraphics%
{gesamttex/edit_VIII-4/images/LH_35_10_14_011_d3_011v.pdf}
\edcaption[0.5em]{%
\lbrack\textit{Fig.\:3}\rbrack%
}%
\end{edfigure}%
\pstart%
\noindent%
Sed si ponamus
fila esse multiplicata in ratione distantiarum,%
\indexS{filum tensum}
ita ut
\textit{A} et \textit{B} connectantur filo uno,
\textit{B} et \textit{C} duobus,
\textit{C} et \textit{D} tribus,
tunc primum vis repraesentanda%
\indexS{vis repraesentanda}
erit
\edtext{%
triangulo,
et quidem \textit{VR} erit%
}{%
\lemma{triangulo,}%
\Bfootnote{%
\textit{(1)}~sit \textit{B} et %
\textit{(2)}~et quidem %
\textit{(a)}~Area \textit{AVR} dabitur ut %
\textit{(b)}~\textit{VR} erit%
~\textit{L}%
}}
vis ascribenda ipsi corpori \textit{V}
soli seu ultimo%
\lbrack,\rbrack\
at $VR + \ind{1}B\ind{1}C$
ascribetur ipsis \textit{V}
et \textit{\ind{1}B},
\edtext{%
\lbrack et\rbrack%
}{%
\lemma{ut}%
\Bfootnote{%
\textit{L\;ändert Hrsg.}%
}}
$VR + \ind{1}B\ind{1}C + \ind{2}B\ind{2}C$
ipsis \textit{V} et \textit{\ind{1}B} et \textit{\ind{2}B},
et ita porro.
Ergo \textit{V} accipiet
\rule[-0mm]{0pt}{5mm}%
$\displaystyle \frac{VR}{1} + \frac{\ind{1}B\ind{1}C}{2} + \frac{\ind{2}B\ind{2}C}{3}$.
\edtext{%
Sit $AB = x$ et $AV = a = VR$.
Fiet vis detrahenda ipsi%
\indexS{vis detrahenda}
$\displaystyle [\ind{1}B] =%
}{%
\lemma{Sit}%
\Bfootnote{%
\textit{(1)}~$AVR = x$ ergo corpus \textit{B} habebit %
\textit{(2)}~\bars{sit \textit{streicht Hrsg.}} $AB = x$ %
\lbrack...\rbrack\ detrahenda ipsi \bars{$\displaystyle \ind{1}B\ind{1}C$ \textit{ändert Hrsg.}} $=$%
~\textit{L}%
}}
\frac{\ind{1}B\ind{1}C}{2} + \frac{\ind{2}B\ind{2}C}{3}$
etc. seu
\rule[-0mm]{0pt}{5.25mm}%
$\displaystyle \int\!\!\! \frac{x}{a - x}dx$.
Quodsi etiam corpora
\textit{V}, \textit{\ind{1}B}, \textit{\ind{2}B}, etc.
sint majora in ratione ipsarum \textit{VR}, \textit{BC}
\rule[-0mm]{0pt}{4.0mm}%
seu quodlibet suo filo alligatum,%
\indexS{filum tensum}%
\indexS{corpus alligatum filo}
tunc idem est calculus virium,%
\indexS{calculus virium}
tantum postea celeritates aliae prodeunt,
pro corporum ratione.
Major adhuc variatio,%
\indexS{variatio}
cum pluribus filis%
\indexS{filum tensum}
ejusdem distantiae%
\lbrack,\rbrack\
diversae magnitudinis corpora alligantur.%
\indexS{corpus alligatum filo}
Caeterum ista applicari possunt corpori
quod %
motum in fluido tenaci%
\indexS{fluidum tenax} %
secum ejus partes rapit%
\lbrack;\rbrack\
ibi enim
pro circulorum distantia
augetur multitudo filorum%
\indexS{multitudo filorum}%
\indexS{filum tensum}
et magnitudo orbium.%
\indexS{orbis}
Sed si perpetuo tenacitas%
\indexS{tenacitas fluidi}
manet,
vis demum omnis peribit.%
\indexS{vis perita}
\ledpb%
Forte naturalis est corporum%
\indexS{cohaesio corporum naturalis}
\edtext{%
cohaesio,
ut si moveatur pars anterior%
}{%
\lemma{cohaesio,}%
\Bfootnote{%
\textit{(1)}~sed %
\textit{(2)}~ut si %
\textit{(a)}~corpus %
\textit{(b)}~moveatur pars anterior%
~\textit{L}%
}}
secum trahet posteriorem nonnihil%
\lbrack,\rbrack\
morantem magis in circulo.
Et
\edtext{%
ita bis non%
}{%
\lemma{ita}%
\Bfootnote{%
\textit{(1)}~semel %
\textit{(2)}~bis%
~\textit{L}%
}}
minuetur,
sed solum distribuetur.
\pend%
\vspace{4.0em}%
\pstart%
\noindent%
\centering%
*     *     *
\pend%
\vspace{4.0em}%
\selectlanguage{german}%
\frenchspacing%
\footnotesize%
\pstart%
\noindent%
\textls{Ausführliche Datierungsgründe:}%
\pagelabel{60117_dat}
Die titellose und undatierte Aufzeichnung N.\;\ref{60117} ist auf einem Träger überliefert, der kein Wasserzeichen aufweist, wohl aber Spuren eines an den in Wien\indexO{Wien} ansässigen kaiserlichen Rat Daniel Nessel\indexN{\textls{Nessel} (Nesselius), Daniel 1644\textendash1700} gerichteten Briefes (vgl. Bl.\;\href{\Scan{LH_35_10_14_011v}}{11\:v\textsuperscript{o}}).
Dieser war seit 1680 Bibliothekar am Habsburgischen Hof und starb dort spätestens am 23.\;März 1700.
Der Briefwechsel zwischen ihm und Leibniz begann im Oktober 1688 in Wien und endete wohl im September 1691.
In diesem Zeitraum entstand wahrscheinlich auch die Aufzeichnung N.\;\ref{60117}.
\pend%
\pstart%
Der besagte Träger %
zeigt jedoch hinsichtlich seiner äußeren Eigenschaften keine Verwandschaft mit den Konzepten oder Abschriften der Briefe an Nessel,\indexN{\textls{Nessel} (Nesselius), Daniel 1644\textendash1700} die Leibniz im Winter 1688/89 aus Wien\indexO{Wien} oder im September 1691 aus Braunschweig\indexO{Braunschweig} abschickte.
Die einzige Ausnahme stellt sein Brief vom 4./14.\;Marz 1691 aus Hannover\indexO{Braunschweig} dar (\textit{LSB} I,\;6 N.\;221).\cite{01484}
Hierbei handelt es sich um Auszüge des Schreibers Gottfried Christian Otto\indexN{\textls{Otto}, Gottfried Christian, geb. 1670, Sterbedatum unbekannt} aus der (heute verschollenen) Anfertigung des benannten Briefes an Nessel,\indexN{\textls{Nessel} (Nesselius), Daniel 1644\textendash1700}.
Dieser Zeuge (LBr~680 Bl.\;13--14) weist hinsichtlich der Farbe, der Textur und der Alterung Merkmale auf, welche denen des Trägers von N.\;\ref{60117} stark ähneln.
Demgemäß erweist sich als plausibel, dass auch N.\;\ref{60117} ungefähr zur gleichen Zeit wie dieser Brief an Nessel\indexN{\textls{Nessel} (Nesselius), Daniel 1644\textendash1700} entstand, d.\,h. im März 1691.
Eine frühere oder spätere Datierung der Aufzeichnung N.\;\ref{60117} ist zwar möglich, wahrscheinlich aber nur innerhalb der Zeitspanne vom Oktober 1688 bis zum September 1691.
\pend%
\selectlanguage{latin}%
\frenchspacing%
\normalsize%
\newpage%
